% Part: history
% Chapter: biographies
% Section: bertrand-russell

\documentclass[../../../include/open-logic-section]{subfiles}

\begin{document}

\olfileid{his}{bio}{rus}

\olsection{বার্ট্রান্ড রাসেল}

\olphoto{russell-bertrand}{বার্ট্রান্ড রাসেল}

বার্ট্রান্ড রাসেলকে আধুনিক বিশ্লেষণী দর্শনের
অন্যতম প্রতিষ্ঠাতা বলা হয়। ১৮৭২ সালের ১৮ মে
জন্ম নেওয়া রাসেল দর্শন ও যুক্তিবিদ্যার কাজের
পাশাপাশি নানা বিষয়ে সাধারণ পাঠকের জন্য
অনেক বই লিখেছিলেন। সারা জীবন তিনি ছিলেন
একজন নিষ্ঠাবান রাজনৈতিক কর্মীও।

ওয়েলসের মনমাউথশায়ারের ট্রেলেখে রাসেলের জন্ম।
তাঁর বাবা-মা ব্রিটিশ অভিজাত শ্রেণির মানুষ
ছিলেন। তাঁরা ছিলেন মুক্তচিন্তক; এমনকি সেই
সময়ে বস্টনের বিদ্রোহী চিন্তাবিদদের সঙ্গেও
তাঁদের বন্ধুত্ব হয়েছিল। দুর্ভাগ্যবশত, শৈশবেই
রাসেলের বাবা-মা মারা যান এবং তাঁকে
দাদু-ঠাকুমার কাছে থাকতে পাঠানো হয়।
সেখানে তাঁকে ধর্মীয় অনুশাসনের মধ্যে বড়
করা হয়—যা তাঁর বাবা-মা যেকোনো মূল্যে
এড়াতে চেয়েছিলেন। নৈতিকতার সব বিষয়ে
তাঁর ঠাকুমা ছিলেন অত্যন্ত কঠোর। কৈশোরে
রাসেল প্রধানত গৃহশিক্ষকদের কাছে বাড়িতেই
পড়াশোনা করেন।

বিশ্লেষণী দর্শনে, বিশেষত যুক্তিবিদ্যায়,
রাসেলের প্রভাব বিরাট। কেমব্রিজের ট্রিনিটি
কলেজে তিনি গণিত ও দর্শন পড়েন; সেখানে
গণিতবিদ ও দার্শনিক আলফ্রেড নর্থ
হোয়াইটহেডের প্রভাব তাঁর উপর পড়ে।
১৯১০ সালে রাসেল ও হোয়াইটহেড
\emph{Principia Mathematica}-র প্রথম খণ্ড
প্রকাশ করেন; এতে তাঁরা গণিতকে যুক্তিবিদ্যায়
নামিয়ে আনা যায়—এই মতের পক্ষে যুক্তি
দেন। পরে রাসেল শত শত বই, প্রবন্ধ ও
রাজনৈতিক পুস্তিকা প্রকাশ করেন। ১৯৫০
সালে তিনি সাহিত্যে নোবেল পুরস্কার পান।

রাজনীতি ও সামাজিক আন্দোলনে রাসেল
গভীরভাবে জড়িত ছিলেন। প্রথম বিশ্বযুদ্ধের
সময়ে যুদ্ধবিরোধী কাজ ও প্রতিবাদের জন্য
তাঁকে গ্রেপ্তার করে ছয় মাস কারাবন্দি রাখা
হয়। কারাগারে তিনি পড়তে ও লিখতে
পেরেছিলেন; তাঁর দাবি, অভিজ্ঞতাটি
“বেশ মনোরম” ছিল। রাসেল ১৯৬১ সালে পারমাণবিক
নিরস্ত্রীকরণের এক সমাবেশে যোগ দেওয়ার
জন্য আবার কারাবন্দি হন। ১৯৪৮ সালে
তিনি একটি বিমান দুর্ঘটনা থেকেও বেঁচে
যান; সেবার কেবল ধূমপানের জন্য
নির্ধারিত অংশের যাত্রীরাই বেঁচেছিলেন।
তাই রাসেল বলতেন, ধূমপানের কাছেই
তাঁর জীবন ঋণী। রাসেলের চারবার বিয়ে
হয়েছিল, তবে বিবাহবহির্ভূত সম্পর্কের জন্যও
তাঁর পরিচিতি ছিল। ১৯৭০ সালের ২
ফেব্রুয়ারি ৯৭ বছর বয়সে ওয়েলসের
পেনরিনদেয়ুদ্রাইথে (Penrhyndeudraeth)
তাঁর মৃত্যু হয়।
% BN-SRC-852 TARGET-CORRECTION-BEGIN
\par\noindent\textbf{পাঠ-সংশোধন।} “সারা জীবন” যুদ্ধবিরোধী থাকার দাবিটি সংশোধিত।
রাসেল-সংরক্ষণাগারের ক্যাটালগে দ্বিতীয় বিশ্বযুদ্ধ নিয়ে ১৯৪০ সালে
তাঁর অবস্থান বদলের কথা আছে; দেখো
\href{https://www.mcmaster.ca/russdocs/Feinberg\%20RA1\%20Catalogue.pdf}{ফাইনবার্গের ক্যাটালগ}, পৃ.~২১৬।
% BN-SRC-852 TARGET-CORRECTION-END

\begin{reading}
১৮৭২--১৯৬৭ সালের জীবন নিয়ে রাসেল
তিন খণ্ডে আত্মজীবনী লিখেছিলেন
\citep{Russell1967,Russell1968,Russell1969}।
ম্যাকমাস্টার বিশ্ববিদ্যালয়ের বার্ট্রান্ড রাসেল
গবেষণাকেন্দ্রে তাঁর আর্কাইভ রয়েছে।
তাঁর সংকলিত রচনার খণ্ডগুলি (অনুসন্ধানযোগ্য
সূচিসহ) এবং আর্কাইভ প্রকল্প সম্পর্কে
জানতে কেন্দ্রটির ওয়েবসাইট দেখো
\citet{Duncan2015}। রাসেলের
\emph{On Denoting} প্রবন্ধটি
\citep{Russell1905} বিশ শতকের
বিশ্লেষণী দর্শনের একটি ধ্রুপদি রচনা।

রাসেল সম্পর্কে স্ট্যানফোর্ড এনসাইক্লোপিডিয়া
অব ফিলসফির নিবন্ধে \citep{Irvine2015}
ইচ্ছা ও রাজনৈতিক তত্ত্ব নিয়ে রাসেলের
কথা বলার অডিও অংশ আছে। তাঁর বহু
ভিডিও সাক্ষাৎকারও অনলাইনে পাওয়া যায়।
যেমন, ধূমপান ও বিমান দুর্ঘটনার অভিজ্ঞতা
নিয়ে তাঁর বক্তব্য দেখতে
\citet{RussellND} দেখো। রাসেলের
\emph{Introduction to Mathematical Philosophy}
সহ কিছু রচনা বিনামূল্যের অডিওবই হিসেবে
\citet{LibriVoxND}-তে পাওয়া যায়।
\end{reading}

\end{document}
